\documentclass[aspectratio=169]{beamer}
\usepackage{iftex}
\ifPDFTeX
  \usepackage[T1]{fontenc}
  \usepackage{lmodern}
\fi
\usepackage[czech]{babel}
\usepackage{amsmath}
\usepackage{xcolor}
\usepackage{listings}
\usepackage{tikz}
\usetikzlibrary{arrows.meta,babel,positioning,shapes.geometric,shapes.misc,shapes.symbols}
\newcommand{\CVUTTemplatePath}{../../../utils/presentation-template}
\usepackage{\CVUTTemplatePath/beamerthemeCVUT}
\setbeamertemplate{footline}[frame number]
\beamertemplatenavigationsymbolsempty

\newcommand{\LANG}{CZ}
\newcommand{\FRONTPAGE}{dark}
\title{Základy návrhu algoritmů}
\subtitle{Základy programování\\Matouš Cejnek\\U12110, FS, ČVUT v Praze}
\author{}
\institute{}
\date{}
\renewcommand{\headlineA}{}
\renewcommand{\headlineB}{}
\renewcommand{\headlineC}{}

\definecolor{CodeBlue}{RGB}{38,83,125}
\definecolor{SoftGray}{RGB}{242,244,247}
\definecolor{DecisionOrange}{RGB}{242,153,74}
\definecolor{SuccessGreen}{RGB}{77,145,95}
\setbeamerfont{normal text}{size=\large}
\setbeamerfont{frametitle}{size=\LARGE}
\setbeamercolor{block title}{fg=white,bg=CodeBlue}
\setbeamercolor{block body}{fg=black,bg=SoftGray}
\setbeamercolor{alertblock title}{fg=white,bg=DecisionOrange}
\setbeamercolor{alertblock body}{fg=black,bg=SoftGray}
\newcommand{\term}[1]{\textcolor{CodeBlue}{\textbf{#1}}}

\lstdefinestyle{python}{
  language=Python,
  basicstyle=\ttfamily\small,
  keywordstyle=\color{CodeBlue}\bfseries,
  commentstyle=\color{SuccessGreen},
  stringstyle=\color{DecisionOrange!80!black},
  showstringspaces=false,
  columns=fullflexible,
  keepspaces=true,
  frame=single,
  rulecolor=\color{black!20},
  backgroundcolor=\color{SoftGray}
}
\tikzset{
  flow/.style={draw=CodeBlue,very thick,fill=SoftGray,align=center,minimum height=9mm},
  process/.style={flow,rectangle,rounded corners=1pt,minimum width=30mm},
  terminator/.style={flow,rounded rectangle,minimum width=25mm},
  decision/.style={flow,diamond,aspect=2.1,inner sep=1pt},
  io/.style={flow,trapezium,trapezium left angle=70,trapezium right angle=110,minimum width=29mm},
  action/.style={flow,rounded corners=4pt,minimum width=30mm},
  arrow/.style={-{Latex[length=2.5mm]},very thick,draw=CodeBlue}
}

\begin{document}

% 0:00--0:01
\begin{frame}
  \thispagestyle{empty}\titlepage
\end{frame}

% 0:01--0:03
\begin{frame}{Co si dnes odnesete?}
  \begin{itemize}
    \item Jak podmínky a cykly řídí \term{tok programu}?
    \item Jak z programu odhadnout \term{počet operací}?
    \item Co říká časová složitost a značení \term{Big O}?
    \item Jak číst a kreslit \term{flowchart}?
    \item Jak souvisí flowchart s \term{UML activity diagramem}?
  \end{itemize}
\end{frame}

\begin{frame}{Na co navazujeme a proč to potřebujeme?}
  \begin{itemize}
    \item Algoritmické myšlení lze rozvíjet bez ohledu na znalost konkrétního programovacího jazyka.
    \item Začneme plánováním řešení \textbf{dříve}, než přejdeme k zápisu kódu.
    \item Stejný problém může mít více správných, ale různě efektivních řešení.
    \item Diagram pomáhá o programu komunikovat bez závislosti na syntaxi jazyka.
  \end{itemize}
\end{frame}

% 0:03--0:07
\begin{frame}{Algoritmus není totéž co program}
  \begin{block}{Algoritmus}
    Konečný a jednoznačný postup, který převádí vstup na požadovaný výstup.
  \end{block}
  \begin{itemize}
    \item \term{Algoritmus} popisuje myšlenku řešení.
    \item \term{Program} je realizace algoritmu v konkrétním jazyce.
    \item Jeden algoritmus lze zapsat slovně, pseudokódem, diagramem i kódem.
  \end{itemize}
\end{frame}

\begin{frame}{Co musí být v zadání problému?}
  \begin{columns}[T]
    \begin{column}{0.46\textwidth}
      \textbf{Vstup}
      \begin{itemize}
        \item Jaká data dostáváme?
        \item Jak jsou velká?
        \item Jaká omezení platí?
      \end{itemize}
    \end{column}
    \begin{column}{0.46\textwidth}
      \textbf{Výstup}
      \begin{itemize}
        \item Co přesně vracíme?
        \item Jak poznáme správnost?
        \item Co se stane na hranicích?
      \end{itemize}
    \end{column}
  \end{columns}
  \begin{alertblock}{Nejdříve specifikace, potom řešení}
    Bez přesného vstupu a výstupu nelze spolehlivě posoudit správnost algoritmu.
  \end{alertblock}
\end{frame}

\begin{frame}{Tři stavební kameny toku programu}
  \centering
  \begin{tikzpicture}[node distance=12mm]
    \node[process] (seq) {posloupnost\\krok za krokem};
    \node[decision,right=of seq] (branch) {podmínka?};
    \node[process,right=of branch] (loop) {opakování\\cyklus};
    \draw[arrow] (seq)--(branch);
    \draw[arrow] (branch)--(loop);
  \end{tikzpicture}
  \vspace{7mm}
  \begin{itemize}
    \item \term{Posloupnost}: příkazy se vykonají v pořadí.
    \item \term{Větvení}: další krok závisí na podmínce.
    \item \term{Cyklus}: část postupu se opakuje.
  \end{itemize}
\end{frame}

% 0:07--0:16
\begin{frame}[fragile]{Podmínka \texttt{if}: vyber jednu cestu}
\begin{lstlisting}[style=python]
if temperature < 0:
    print("Mrzne")
else:
    print("Nemrzne")
\end{lstlisting}
  \begin{itemize}
    \item Podmínka má hodnotu \texttt{True} nebo \texttt{False}.
    \item Vykoná se právě odpovídající větev.
    \item Po větvení program pokračuje společným dalším krokem.
  \end{itemize}
\end{frame}

\begin{frame}[fragile]{Více možností: \texttt{if} -- \texttt{elif} -- \texttt{else}}
\begin{lstlisting}[style=python]
if score >= 90:
    grade = "A"
elif score >= 75:
    grade = "B"
else:
    grade = "C"
\end{lstlisting}
  \begin{alertblock}{Pořadí podmínek je součást algoritmu}
    Python testuje podmínky shora dolů a použije první pravdivou větev.
  \end{alertblock}
\end{frame}

\begin{frame}[fragile]{Cyklus \texttt{while}: opakuj, dokud podmínka platí}
\begin{lstlisting}[style=python]
n = 8
while n > 1:
    n = n // 2
    print(n)
\end{lstlisting}
  \begin{itemize}
    \item Podmínka se kontroluje \textbf{před každým průchodem}.
    \item Tělo se nemusí vykonat ani jednou.
    \item Některý krok musí směřovat k ukončení cyklu.
  \end{itemize}
\end{frame}

\begin{frame}{Kdy se \texttt{while} nezastaví?}
  \begin{itemize}
    \item Řídicí hodnota se vůbec nemění.
    \item Mění se opačným směrem, než potřebuje podmínka.
    \item Některá větev přeskočí aktualizaci hodnoty.
    \item Ukončení závisí na vstupu, který nemusí nikdy nastat.
  \end{itemize}
  \begin{block}{Invariant a varianta cyklu}
    \term{Invariant} zůstává pravdivý; \term{varianta} se mění směrem k ukončení.
  \end{block}
\end{frame}

\begin{frame}[fragile]{Cyklus \texttt{for}: projdi známou kolekci}
\begin{lstlisting}[style=python]
numbers = [3, 8, 1, 6]
total = 0

for number in numbers:
    total += number
\end{lstlisting}
  \begin{itemize}
    \item Proměnná \texttt{number} postupně získá každý prvek.
    \item Počet průchodů je dán počtem prvků kolekce.
    \item V Pythonu \texttt{for} znamená „pro každý prvek“.
  \end{itemize}
\end{frame}

\begin{frame}{\texttt{for}, nebo \texttt{while}?}
  \begin{columns}[T]
    \begin{column}{0.46\textwidth}
      \textbf{Použijte \texttt{for}}
      \begin{itemize}
        \item pro každý prvek,
        \item známý počet opakování,
        \item průchod přes \texttt{range}.
      \end{itemize}
    \end{column}
    \begin{column}{0.46\textwidth}
      \textbf{Použijte \texttt{while}}
      \begin{itemize}
        \item opakování do splnění podmínky,
        \item počet kroků předem neznáme,
        \item stav řídíme sami.
      \end{itemize}
    \end{column}
  \end{columns}
\end{frame}

\begin{frame}[fragile]{\texttt{break} a \texttt{continue} mění běžný tok}
\begin{lstlisting}[style=python]
for number in numbers:
    if number < 0:
        continue       # dalsi prvek
    if number == 0:
        break          # konec cyklu
    total += number
\end{lstlisting}
  \begin{itemize}
    \item \texttt{continue} přeskočí zbytek aktuálního průchodu.
    \item \texttt{break} ukončí nejbližší cyklus.
    \item V diagramu musí být obě změny toku viditelné šipkou.
  \end{itemize}
\end{frame}

% 0:16--0:31
\begin{frame}{Proč počítat operace?}
  \begin{itemize}
    \item Čas na stopkách závisí na počítači, jazyce i aktuálním zatížení.
    \item Počet kroků algoritmu lze porovnávat nezávisle na konkrétním stroji.
    \item Zajímá nás, jak práce roste s velikostí vstupu.
  \end{itemize}
  \begin{block}{Velikost vstupu \(n\)}
    Například počet prvků seznamu, počet znaků textu nebo počet vrcholů grafu.
  \end{block}
\end{frame}

\begin{frame}{Co budeme považovat za jednu operaci?}
  \begin{itemize}
    \item přiřazení nebo aritmetická operace,
    \item porovnání a vyhodnocení jednoduché podmínky,
    \item přístup k jednomu prvku seznamu,
    \item jednoduchý výpis --- v našem zjednodušeném modelu.
  \end{itemize}
  \begin{alertblock}{Je to model}
    Skutečné operace netrvají stejně dlouho. Model musí být jednoduchý, ale použitý konzistentně.
  \end{alertblock}
\end{frame}

\begin{frame}[fragile]{Příklad: přesný počet kroků}
\begin{lstlisting}[style=python]
total = 0                    # 1x
for number in numbers:       # n prvku
    total = total + number   # n-krat
print(total)                 # 1x
\end{lstlisting}
  \begin{itemize}
    \item Zjednodušeně: \(T(n)=1+n+n+1=2n+2\).
    \item Přesný výsledek závisí na tom, co zahrneme do jedné operace.
    \item Důležitější než koeficient je zde \textbf{lineární růst}.
  \end{itemize}
\end{frame}

\begin{frame}{Podmínka: kterou větev počítáme?}
  \begin{itemize}
    \item \term{Nejlepší případ}: vstup vyvolá nejméně práce.
    \item \term{Nejhorší případ}: vstup vyvolá nejvíce práce.
    \item \term{Průměrný případ}: potřebujeme předpoklad o rozdělení vstupů.
  \end{itemize}
  \begin{block}{Běžná dohoda}
    Není-li uvedeno jinak, často analyzujeme nejhorší případ: dává bezpečnou horní mez.
  \end{block}
\end{frame}

\begin{frame}[fragile]{Předčasný konec: nejlepší a nejhorší případ}
\begin{lstlisting}[style=python]
for number in numbers:
    if number == target:
        return True
return False
\end{lstlisting}
  \begin{itemize}
    \item Cíl je první prvek: přibližně \(1\) kontrola.
    \item Cíl je poslední nebo chybí: přibližně \(n\) kontrol.
    \item Nejlepší případ \(\Theta(1)\), nejhorší \(\Theta(n)\).
  \end{itemize}
\end{frame}

\begin{frame}[fragile]{Dva cykly za sebou se sčítají}
\begin{lstlisting}[style=python]
for x in numbers:       # n
    print(x)

for x in numbers:       # n
    print(x * x)
\end{lstlisting}
  \[
    T(n) \approx n+n=2n \quad\Rightarrow\quad \Theta(n)
  \]
  Dvojnásobek práce je stále lineární růst.
\end{frame}

\begin{frame}[fragile]{Vnořené cykly se násobí}
\begin{lstlisting}[style=python]
for x in numbers:          # n
    for y in numbers:      # n pro kazde x
        print(x, y)
\end{lstlisting}
  \[
    T(n) \approx n\cdot n=n^2 \quad\Rightarrow\quad \Theta(n^2)
  \]
  Každá z \(n\) vnějších iterací spustí celý vnitřní cyklus.
\end{frame}

\begin{frame}[fragile]{Půlení problému vede k logaritmu}
\begin{lstlisting}[style=python]
while n > 1:
    n = n // 2
\end{lstlisting}
  Po \(k\) krocích zbývá přibližně \(n/2^k\). Končíme, když
  \[
    \frac{n}{2^k}\leq 1 \quad\Rightarrow\quad k\geq \log_2 n.
  \]
  Proto je počet průchodů \(\Theta(\log n)\).
\end{frame}

\begin{frame}{Od přesného počtu k řádu růstu}
  \begin{itemize}
    \item \(3n+7\), \(100n+2\) i \(n\) rostou pro velké \(n\) lineárně.
    \item Konstantní násobky a nižší členy nemění dominantní růst.
    \item Asymptotická analýza se ptá: \textbf{co převládne pro velké vstupy?}
  \end{itemize}
  \[
    4n^2+3n+12 \quad\leadsto\quad n^2
  \]
\end{frame}

\begin{frame}{Big O: asymptotická horní mez}
  \begin{block}{Definice}
    \(T(n)\in O(g(n))\), pokud existují konstanty \(c>0\) a \(n_0\), že pro všechna \(n\geq n_0\) platí \(T(n)\leq c\,g(n)\).
  \end{block}
  \begin{itemize}
    \item Big O říká, že růst \(T\) není asymptoticky rychlejší než \(g\).
    \item Formálně je \(3n+2\) také v \(O(n^2)\), ale těsnější mez je \(O(n)\).
    \item V praxi se snažíme uvádět co nejtěsnější užitečnou mez.
  \end{itemize}
\end{frame}

\begin{frame}{Big O, Omega a Theta}
  \begin{itemize}
    \item \(O(g(n))\): asymptotická \textbf{horní} mez.
    \item \(\Omega(g(n))\): asymptotická \textbf{dolní} mez.
    \item \(\Theta(g(n))\): horní i dolní mez --- stejný řád růstu.
  \end{itemize}
  \begin{block}{Příklad}
    Pro \(T(n)=3n+2\) platí \(T(n)\in O(n)\), \(T(n)\in\Omega(n)\), tedy \(T(n)\in\Theta(n)\).
  \end{block}
\end{frame}

\begin{frame}{Typické třídy růstu}
  \centering
  \begin{tabular}{lll}
    \textbf{Složitost} & \textbf{Název} & \textbf{Typický příklad} \\
    \hline
    \(O(1)\) & konstantní & přístup na známý index \\
    \(O(\log n)\) & logaritmická & opakované půlení \\
    \(O(n)\) & lineární & jeden průchod seznamem \\
    \(O(n\log n)\) & linearitmická & efektivní třídění \\
    \(O(n^2)\) & kvadratická & všechny dvojice prvků \\
    \(O(2^n)\) & exponenciální & výběr všech podmnožin \\
  \end{tabular}
\end{frame}

\begin{frame}{Malý vstup může klamat}
  \begin{itemize}
    \item Algoritmus s \(1000n\) může být pro malé \(n\) pomalejší než algoritmus s \(n^2\).
    \item Pro dost velké \(n\) však kvadratický růst lineární překoná.
    \item Big O neříká skutečný čas ani konstanty.
    \item Po asymptotické analýze má smysl program také měřit.
  \end{itemize}
\end{frame}

% 0:31--0:50
\begin{frame}{Co je flowchart?}
  \begin{block}{Vývojový diagram (flowchart)}
    Grafické znázornění kroků algoritmu a toku řízení mezi nimi.
  \end{block}
  \begin{itemize}
    \item Uzly představují kroky, rozhodnutí nebo vstup a výstup.
    \item Orientované šipky určují pořadí vykonávání.
    \item Diagram zachycuje logiku; nemusí kopírovat každý řádek kódu.
  \end{itemize}
\end{frame}

\begin{frame}{Mermaid je zápis, flowchart je diagram}
  \begin{itemize}
    \item \term{Flowchart} je způsob modelování algoritmu.
    \item \term{Mermaid} je textový jazyk, který diagram vykreslí.
    \item Tentýž flowchart lze nakreslit rukou, v editoru nebo zapsat v Mermaid.
  \end{itemize}
  \begin{block}{Příklad zápisu}
    \texttt{flowchart TD}\qquad \texttt{start([Start]) --> process[Zpracuj data]}
  \end{block}
\end{frame}

\begin{frame}{Značka: začátek a konec}
  \centering
  \begin{tikzpicture}[node distance=35mm]
    \node[terminator] (start) {Start};
    \node[terminator,right=of start] (stop) {Konec};
    \draw[arrow] (start)--(stop);
  \end{tikzpicture}
  \vspace{8mm}
  \begin{itemize}
    \item Zaoblený terminátor označuje vstupní nebo koncový bod toku.
    \item Diagram má mít jasný začátek a alespoň jeden konec.
    \item Mermaid: \texttt{start([Start])}, \texttt{stop([Konec])}.
  \end{itemize}
\end{frame}

\begin{frame}{Značka: proces neboli akce}
  \centering
  \begin{tikzpicture}
    \node[process,minimum width=50mm] {\texttt{total = total + number}};
  \end{tikzpicture}
  \vspace{8mm}
  \begin{itemize}
    \item Obdélník představuje výpočet, přiřazení nebo jiný krok.
    \item Text má být krátký, konkrétní a ve slovesném tvaru nebo jako příkaz.
    \item Mermaid: \texttt{step[Sečti hodnoty]}.
  \end{itemize}
\end{frame}

\begin{frame}{Značka: rozhodnutí}
  \centering
  \begin{tikzpicture}[node distance=18mm,scale=.78,transform shape]
    \node[decision] (d) {\texttt{n > 0?}};
    \node[process,above right=of d] (yes) {kladné};
    \node[process,below right=of d] (no) {nekladné};
    \draw[arrow] (d)--node[above left]{ano}(yes);
    \draw[arrow] (d)--node[below left]{ne}(no);
  \end{tikzpicture}
  \vspace{1mm}
  \begin{itemize}
    \item Kosočtverec obsahuje otázku nebo logickou podmínku.
    \item Každá odchozí hrana musí mít jednoznačný popisek.
    \item Mermaid: \texttt{decision\{n > 0?\}}.
  \end{itemize}
\end{frame}

\begin{frame}{Značka: vstup a výstup}
  \centering
  \begin{tikzpicture}[node distance=28mm]
    \node[io] (input) {Načti \texttt{n}};
    \node[io,right=of input] (output) {Vypiš výsledek};
    \draw[arrow] (input)--(output);
  \end{tikzpicture}
  \vspace{8mm}
  \begin{itemize}
    \item Rovnoběžník tradičně označuje načtení nebo výpis dat.
    \item Odlišuje komunikaci s okolím od vnitřního výpočtu.
    \item Mermaid podporuje tvar rovnoběžníku; běžný obdélník je také čitelný, pokud jsme konzistentní.
  \end{itemize}
\end{frame}

\begin{frame}{Značka: šipka a spojení toku}
  \centering
  \vspace{4mm}
  \begin{tikzpicture}[node distance=22mm]
    \node[process] (a) {krok A};
    \node[process,right=of a] (b) {krok B};
    \node[process,right=of b] (c) {krok C};
    \draw[arrow] (a)--(b);
    \draw[arrow] (b)--(c);
  \end{tikzpicture}
  \vspace{8mm}
  \begin{itemize}
    \item Hrot šipky určuje směr řízení.
    \item Tok obvykle vedeme shora dolů nebo zleva doprava.
    \item Křížení hran omezujeme; složitý obraz často signalizuje příliš mnoho detailů.
  \end{itemize}
\end{frame}

\begin{frame}{Jak diagram vyjádří cyklus?}
  \centering
  \vspace{4mm}
  \begin{tikzpicture}[node distance=12mm]
    \node[decision] (condition) {další prvek?};
    \node[process,below=of condition] (body) {zpracuj prvek};
    \node[terminator,right=32mm of condition] (end) {pokračuj dál};
    \draw[arrow] (condition)--node[right]{ano}(body);
    \draw[arrow] (body.west) to[out=180,in=180,looseness=1.4] (condition.west);
    \draw[arrow] (condition)--node[above]{ne}(end);
  \end{tikzpicture}
  \vspace{4mm}
  \begin{itemize}
    \item Rozhodnutí testuje pokračování cyklu.
    \item Zpětná šipka vrací tok k další kontrole podmínky.
  \end{itemize}
\end{frame}

\begin{frame}{Další tvary: užitečné, ne vždy nutné}
  \centering
  \begin{tikzpicture}[node distance=19mm]
    \node[flow,cylinder,shape border rotate=90,aspect=.3,minimum width=25mm] (db) {data};
    \node[flow,rectangle,double,double distance=1.5pt,
      minimum width=34mm,right=of db] (sub) {podprogram};
    \node[flow,tape,tape bend bottom=out,tape bend top=none,
      minimum width=28mm,right=of sub] (doc) {dokument};
  \end{tikzpicture}
  \vspace{7mm}
  \begin{itemize}
    \item Válec: uložená data nebo databáze.
    \item Obdélník s dvojitým okrajem: předem definovaný proces / podprogram.
    \item Dokument: dokument nebo report.
    \item Pro jednoduché algoritmy obvykle stačí terminátor, proces, rozhodnutí, I/O a šipky.
  \end{itemize}
\end{frame}

\begin{frame}{Flowchart versus UML activity diagram}
  \begin{columns}[T]
    \begin{column}{0.46\textwidth}
      \textbf{Flowchart}
      \begin{itemize}
        \item obecná a pružná notace,
        \item tradiční tvary pro proces a I/O,
        \item vhodný pro vysvětlení algoritmu.
      \end{itemize}
    \end{column}
    \begin{column}{0.46\textwidth}
      \textbf{UML activity diagram}
      \begin{itemize}
        \item standardizovaná UML notace,
        \item akce, rozhodnutí a řídicí toky,
        \item umí paralelní větve a odpovědnosti.
      \end{itemize}
    \end{column}
  \end{columns}
  \begin{block}{Pro naše úlohy}
    Význam je velmi podobný. Důležitější než název je čitelný a konzistentní tok.
  \end{block}
\end{frame}

\begin{frame}{Jak se liší základní značky UML aktivit?}
  \centering
  \begin{tikzpicture}[node distance=17mm]
    \node[circle,fill=black,minimum size=6mm] (initial) {};
    \node[action,right=of initial] (action) {zpracuj vstup};
    \node[decision,right=of action] (decision) {platí?};
    \node[circle,draw=black,very thick,minimum size=9mm,right=of decision] (outer) {};
    \node[circle,fill=black,minimum size=5mm] at (outer.center) {};
    \draw[arrow] (initial)--(action);
    \draw[arrow] (action)--(decision);
    \draw[arrow] (decision)--node[above]{[ne]}(outer);
  \end{tikzpicture}
  \vspace{7mm}
  \begin{itemize}
    \item Plný kruh: počáteční uzel; terč: konec aktivity.
    \item Zaoblený obdélník: akce; kosočtverec: rozhodnutí nebo sloučení.
    \item Podmínky na hranách se často zapisují jako \texttt{[ano]} a \texttt{[ne]}.
  \end{itemize}
\end{frame}

% 0:50--0:57
\begin{frame}[fragile]{Od programu k diagramu: nejdříve bloky}
  \vspace{3mm}
  \begin{columns}[T]
    \begin{column}{0.47\textwidth}
\begin{lstlisting}[style=python,basicstyle=\ttfamily\footnotesize]
total = 0
for n in numbers:
    if n % 2 == 0:
        total += n
print(total)
\end{lstlisting}
    \end{column}
    \begin{column}{0.47\textwidth}
      \begin{enumerate}
        \item inicializace,
        \item kontrola dalšího prvku,
        \item rozhodnutí sudé / liché,
        \item případná aktualizace,
        \item výpis po skončení cyklu.
      \end{enumerate}
    \end{column}
  \end{columns}
\end{frame}

\begin{frame}{Od programu k diagramu: výsledný tok}
  \centering
  \vspace{4mm}
  \begin{tikzpicture}[node distance=4mm,scale=.75,transform shape]
    \node[terminator] (start) {Start};
    \node[process,below=of start] (init) {\texttt{total = 0}};
    \node[decision,below=of init] (more) {další \texttt{n}?};
    \node[decision,below=of more] (even) {\texttt{n \% 2 == 0?}};
    \node[process,below=of even] (add) {\texttt{total = total + n}};
    \node[io,right=32mm of more] (out) {vypiš \texttt{total}};
    \node[terminator,below=of out] (stop) {Konec};
    \draw[arrow] (start)--(init);
    \draw[arrow] (init)--(more);
    \draw[arrow] (more)--node[right]{ano}(even);
    \draw[arrow] (even)--node[right]{ano}(add);
    \draw[arrow] (add.west) to[out=180,in=180,looseness=1.3] (more.west);
    \draw[arrow] (even.east) to[out=0,in=0,looseness=1.5]
      node[pos=.25,right]{ne} (more.east);
    \draw[arrow] (more)--node[above]{ne}(out);
    \draw[arrow] (out)--(stop);
  \end{tikzpicture}
\end{frame}

\begin{frame}{Časté chyby v diagramech}
  \begin{itemize}
    \item Rozhodnutí nemá popsané odchozí větve.
    \item Cyklus nemá zpětnou hranu nebo cestu ven.
    \item Výpis uvnitř cyklu je omylem zakreslen až za cyklem.
    \item \texttt{break} vede k další iteraci místo za cyklus.
    \item Diagram mechanicky kopíruje každý řádek a přestává být čitelný.
  \end{itemize}
\end{frame}

\begin{frame}{Jak kontrolovat správnost diagramu?}
  \begin{enumerate}
    \item Zvolte malý konkrétní vstup.
    \item Sledujte šipky a zapisujte hodnoty proměnných.
    \item Stejný vstup projděte řádek po řádku v programu.
    \item Porovnejte pořadí kroků, větve a výsledný výstup.
    \item Zopakujte pro hraniční vstup: prázdný seznam, nulu, první shodu.
  \end{enumerate}
\end{frame}

% 0:57--1:00
\begin{frame}{Jedna úloha, tři pohledy}
  \begin{itemize}
    \item \term{Kód} přesně říká, co vykoná počítač.
    \item \term{Diagram} ukazuje strukturu toku a usnadňuje komunikaci.
    \item \term{Analýza složitosti} říká, jak množství práce roste se vstupem.
  \end{itemize}
  \begin{block}{Dobrý návrh propojuje všechny tři}
    Dokážeme vysvětlit správnost, zakreslit tok a odhadnout počet operací.
  \end{block}
\end{frame}

\begin{frame}{Co si zapamatovat?}
  \begin{itemize}
    \item \texttt{if} větví tok; \texttt{for} prochází prvky; \texttt{while} opakuje podle podmínky.
    \item Sekvenční cykly sčítáme, vnořené cykly typicky násobíme.
    \item Big O je asymptotická horní mez; \(\Theta\) vyjadřuje těsný řád růstu.
    \item Flowchart tvoří uzly a orientované šipky; značky používejte konzistentně.
    \item UML activity diagram je formálnější příbuzná notace, nikoli jiný algoritmus.
  \end{itemize}
  \vspace{3mm}
  \centering\Large Otázky?
\end{frame}

\end{document}
